% s7-identification.tex -- Section 7 of the reader's companion (drafter: aqft, R3; edited in R4).
% Displays transcribed from part_i/main.tex carry \tag*{\rec{...}{...}}.
\section{From Cauchy to the theorem}\label{sec:identification}

\emph{The ground-state property costs an identity, not an estimate.}

Everything in this section is \condtext{conditional on $\UGJS$} through Theorem~11.1,
except Proposition~12.1 and Lemma~2.8, which do not use A8. Theorem~11.1 turns the
one-path estimate of \S\ref{sec:lppl} into convergence \rec{Theorem~11.1}{28}. Step~1:
if $\Psi_{m_1}(d)<\varepsilon$, all $g,g'\in\G$ whose plateau interiors contain $O_d$
satisfy $\|(\omega_g-\omega_{g'})|_{\A(O)}\|\le\Psi_{m_1}(d)<\varepsilon$ \rec{(32)}{28},
so the net is Cauchy in each $\A(O)^*$. Step~2: completeness of each $\A(O)^*$ gives
limits, compatible under restriction, that define a state $\omega_\infty$ on $\A$.
Step~3: the cutoffs equal to $1$ on $O_d$ are cofinal in $\G$, so the bound passes to
the limit; this is the checkpoint (31) displayed at the end of \S\ref{sec:architecture}
\rec{(31)}{28}.

Local normality comes with the topology: Proposition~11.2 obtains it from the
normality of the vector states $\omega_g$ on the von Neumann algebra $\A(O)$
\rec{Proposition~11.2}{29}.

The dynamics $\alpha$ is imported: A6(ii) supplies a one-parameter automorphism group of
$\A$, independent of $g$, with $\alpha_t(\A(O))\subseteq\A(O_{|t|})$ and the patching
property (P): $\alpha_t(A)=\beta^g_t(A)$ for $A\in\A(O)$ whenever $g\equiv1$ on
$O_{|t|}$ \rec{A6}{9}. No continuity of $\alpha$ is imported or used. If $|t|<d(g,O)$,
then $g\equiv1$ on $O_{|t|}$, and (P) with $e^{-itH(g)}\Om_g=\Om_g$ gives, for
$A,B\in\A(O)$, the identity of Fig.~\ref{fig:window}:
\begin{equation*}
  \omega_g\bigl(A\alpha_t(B)\bigr)
  =\ip{\Om_g}{Ae^{itH(g)}Be^{-itH(g)}\Om_g}
  =\ip{A^*\Om_g}{e^{itH(g)}B\Om_g}=G_g(t) .
  \tag*{\rec{proof of Proposition~12.1}{29}}
\end{equation*}
Here $G_g(t):=\int e^{itE}\,d\rho_g(E)$, where, with $P_{H(g)}$ the spectral measure of
$H(g)$, $\rho_g(\cdot):=\ip{A^*\Om_g}{P_{H(g)}(\cdot)B\Om_g}$ is a complex measure with
$\supp\rho_g\subseteq\spec H(g)\subseteq[0,\infty)$: inside the window, each two-point
function is the Fourier--Stieltjes transform of a measure on $[0,\infty)$. For
$|t|\le T$, $A\alpha_t(B)\in\A(O_T)$ has norm at most $\|A\|\|B\|$, so norm convergence
on $\A(O_T)$ makes these functions converge uniformly on $|t|\le T$, and the definition
passes to the limit:
\begin{recordbox}{Proposition 12.1 of the record (p.\,29)}
If $(g_n)\subseteq\G$ satisfies $d(g_n,O)\to\infty$ for every bounded interval $O$, and
$\omega_{g_n}\to\omega$ in norm on every $\A(O)$, then $\omega$ is a ground state of
$(\A,\alpha)$ in the spectral sense: $\omega\circ\alpha_t=\omega$ for all $t$; the
function $t\mapsto\omega(A\alpha_t(B))$ is continuous and bounded for all $A,B\in\A$; and
\[
 \int_\R f(t)\,\omega(A\alpha_t(B))\,dt=0
\]
for every $f\in L^1(\R)$ whose transform $\check f(E)=\int f(t)e^{itE}\,dt$ vanishes on
$[0,\infty)$.
\end{recordbox}

Translation invariance needs no averaging: $g\mapsto\tau_ag$ is an order isomorphism of
$\G$ with cofinal range, so (T), $\omega_{\tau_ag}=\omega_g\circ\tau_{-a}$
\rec{Lemma~2.8}{12}, and Theorem~12.2(1) give $\omega_\infty\circ\tau_a=\omega_\infty$;
parity follows likewise for even $\mathscr P$ \rec{proof of Theorem~12.2}{30}.

By Corollary~12.3, every $w^*$ limit point of the averaged or unaveraged cutoff states of
\cite[\S2]{GJ3} coincides with $\omega_\infty$ on each $\A(O)$ \rec{Corollary~12.3}{31}.
Its proof takes GJ3's cutoff function $g$ in $\G$ and applies Theorem~12.2(1), uniformly
in the shift, to each shifted cutoff in the average. So $\omega_\infty$ is canonical for
the cutoff family; its uniqueness is not claimed (\S\ref{sec:notclaimed}).

\noindent\emph{In the record:} A6, p.~9; Lemma~2.8, p.~12; Theorem~11.1 with
eqs.~(31)--(32), Propositions~11.2 and 12.1, Theorem~12.2 and Corollary~12.3, pp.~28--31.
